\documentclass[10pt,twocolumn,a4paper]{article}

\usepackage[margin=0.75in]{geometry}
\usepackage{times}
\usepackage{enumitem}
\usepackage{amsmath}
\usepackage{booktabs}
\usepackage{multirow}
\usepackage{xcolor}
\usepackage{listings}
\usepackage[numbers,sort&compress]{natbib}
\usepackage{hyperref}
\usepackage{microtype}
\usepackage{float}
\usepackage{cuted}
\usepackage[font={small,it},labelfont={bf,up}]{caption}

\setlength{\parskip}{5pt}
\setlength{\parindent}{1.5em}
\setlength{\textfloatsep}{14pt plus 2pt minus 2pt}
\setlength{\floatsep}{12pt plus 2pt minus 2pt}
\setlength{\intextsep}{10pt plus 2pt minus 2pt}

% ---- JSON listing style (used in the axiom-validation section) ----
\definecolor{jsongray}{gray}{0.45}
\definecolor{jsonstr}{rgb}{0.10,0.25,0.60}
\definecolor{jsonframe}{gray}{0.75}
\lstdefinestyle{json}{
  basicstyle=\ttfamily\footnotesize,
  breaklines=true,
  breakatwhitespace=true,
  showstringspaces=false,
  columns=fullflexible,
  keepspaces=true,
  frame=single,
  rulecolor=\color{jsonframe},
  framesep=4pt,
  xleftmargin=4pt,
  xrightmargin=2pt,
  stringstyle=\color{jsonstr},
  commentstyle=\color{jsongray}\itshape,
}

\title{\textbf{Aggregating Ethics:}\\ Moral Uncertainty Aggregation for LLM Alignment}
\author{MSc Dissertation Draft}
\date{}

\begin{document}
\raggedbottom
\maketitle

\section{Introduction}

As AI systems become increasingly capable and influential, the problem of
ensuring they behave in accordance with human values grows correspondingly
urgent. The alignment problem \cite{christian2020} is the task of ensuring
that an AI system pursues the goals and honours the values its developers
and users intend, rather than proxies that merely correlate with them. In
practice, the dominant approach to aligning large language models is
Reinforcement Learning from Human Feedback (RLHF) \cite{ouyang2022}, in
which human raters compare model outputs and their preferences are
distilled into a reward signal that guides fine-tuning.

These raters often disagree with one another, particularly in normatively
ambiguous moral dilemmas, where there is no universal consensus on right or
wrong. Crucially, these disagreements do not stem merely from personal
idiosyncrasy or noise. They reflect principled differences between
competing ethical frameworks. A utilitarian, a Kantian, and an Ubuntu
thinker may each reach a different conclusion while reasoning coherently
from within their own tradition. In RLHF, this disagreement is not
preserved, however. A single reward model is trained to fit the aggregate
of many annotators' conflicting preferences, converging on a signal that
reflects the pool's majority tendency and offers no explicit account of how
competing moral views ought to be combined \cite{casper2023}.

\textbf{This dissertation investigates whether the choice of moral
aggregation method produces measurably different LLM behaviour when applied
to normatively ambiguous dilemmas.} Drawing on \citeauthor{macaskill2020}'s
formal framework for decision-making under moral uncertainty
\cite{macaskill2020}, I implement three aggregation schemes (Expected
Choiceworthiness, Maximin, and Nash Moral Parliament) as competing
response-selection criteria for a language model, alongside a single-theory
utilitarian baseline. Candidate responses are scored by three ethical
reward models operationalising three distinct moral traditions, namely
utilitarian, deontological, and Ubuntu. I test whether different
aggregation methods select systematically different responses, and if so,
where and why the methods diverge. This question has growing practical
importance, and as users increasingly turn to language models for guidance on
sensitive personal and ethical decisions, the implicit moral character of
these systems is no longer a theoretical concern. It is a design choice with
real consequences.

Stated precisely, the question is this: \emph{holding fixed a pool of
candidate responses to a morally contested dilemma, and holding fixed a set
of validated judges representing utilitarian, deontological, and Ubuntu
evaluation, does the choice of aggregation rule change which response is
selected, and does it change the action that response recommends?} The two
halves are deliberately distinct. Two rules may prefer different texts that
steer a reader toward the same course of action, which is disagreement over
justification. Or, they may prefer texts whose recommendations are opposed,
which is disagreement over what ought to be done. The design measures both
(Section~\ref{sec:divergence}).

The dissertation makes three contributions.
\begin{enumerate}[leftmargin=*, itemsep=2pt, topsep=2pt]
  \item \textbf{An end-to-end pipeline for moral-uncertainty aggregation.}
  Five hundred everyday dilemmas filtered for cross-framework tension, each
  with sixteen candidate responses. Three framework judges score all
  $8{,}000$ responses, and four selection rules are applied to the
  normalised scores, with Ubuntu included on equal terms with the two
  Western traditions (Section~\ref{sec:method}).
  \item \textbf{A behavioural validation method for framework judges.} Each
  judge is tested against fifteen operational axioms drawn from the
  normative literature, with directional predictions for the other two
  judges and a held-out set guarding against prompt overfitting
  (Section~\ref{sec:axiom-validation}). The harness is reusable for
  certifying any purported framework judge.
  \item \textbf{A two-level, tie-robust measure of divergence.} Selections
  are treated as winner sets rather than broken by arbitrary tie-breaks,
  and disagreement over the chosen text is reported separately from
  disagreement over the recommended action, so that no reported
  disagreement could be removed by a choice of tie-break
  (Section~\ref{sec:divergence}).
\end{enumerate}

Section~\ref{sec:background} sets out the formal problem and the
aggregation proposals this study operationalises,
Section~\ref{sec:related} situates the work in the surrounding literature,
and Sections~\ref{sec:method} and~\ref{sec:results} describe the pipeline
and its results.

\section{Moral Uncertainty}
\label{sec:background}

Moral uncertainty is uncertainty about which moral theory is correct, as
distinct from uncertainty about the facts. An agent may know exactly what an
action will cause, yet still not know whether outcomes are all that morally
matters. Formally, the agent faces options $A$ and entertains rival theories
$T_1, \dots, T_n$ with credences $w_1, \dots, w_n$ summing to one, each
grading the options with a real-valued choiceworthiness function $c_i$. The
problem is which option is rational to choose, given credences and
verdicts together \cite{macaskill2020}. The sceptical reply, "act on the
theory you find most plausible", ignores stakes in the same way that acting
on the most probable state of the world, whatever the payoffs, ignores
stakes under empirical uncertainty. A theory held at credence $0.4$ that
mildly favours an option should not silence theories with combined credence
$0.6$ that condemn it as catastrophic \cite{macaskill2020}.

The standard constructive proposal, maximising \textbf{expected
choiceworthiness} (MEC), imports the apparatus of expected utility
\cite{lockhart2000,macaskill2020}:
\[
  a^{\star} \;=\; \arg\max_{a \in A}\; \sum_{i=1}^{n} w_i\, c_i(a).
\]
MEC presupposes that the $c_i$ are cardinal and expressed in a common unit,
though the theories themselves supply neither. There is no theory-neutral answer
to how many units of utilitarian welfare a broken Kantian duty is worth.
This is the \emph{intertheoretic value problem}, the central obstacle to
any additive rule \cite{macaskill2020}. The usual response is
normalisation, which attempts to place the theories on a common scale by statistical fiat.
\citet{lockhart2000} proposes equalising each theory's range over the
options, and \citet{cottonbarratt2020} argue that equalising variance better
gives every theory an equal say. A further dispute concerns scope. Do we
normalise over the present decision's options alone, or over all options a
theory could ever grade? Both problem and responses recur here, as the judges
differ in how much of the $0$--$10$ scale they use, and the global
$z$-scoring of Section~\ref{sec:zscore} is a broad-scope variance
normalisation.

MEC also inherits the pathologies of expected utility, chief among them
\emph{fanaticism}. Because credence and stakes are multiplied, a theory
held at very low credence can still dominate the sum whenever it declares
the stakes extreme enough, and the agent ends up governed by a theory it
barely believes \cite{macaskill2020}. That worry motivates rules that do
not add. \citet{bostrom2009} pictures the theories as a \textbf{moral
parliament}, delegates who hold seats in proportion to credence and
negotiate a compromise rather than being summed, a proposal later
formalised by \citet{newberry2021}. The parliament is a picture rather
than a rule, and bargaining theory supplies the missing mechanism. The
\textbf{Nash solution} \cite{nash1950}, applied to moral uncertainty by
\citet{greaves2019}, selects the option that maximises the product of
each theory's gain over a disagreement point. This rewards compromise,
since an option that leaves any one theory near its disagreement point
has its product destroyed however well the others fare. Treating the
theories as voters instead imports the impossibility results of social
choice \cite{arrow1951} and, being rank-based, discards the magnitudes
MEC trades on. Finally, a risk-averse agent may prefer to protect the
worst-served theory rather than maximise an average. This yields
\textbf{maximin}, which grades each option by the lowest choiceworthiness
any theory
assigns it, and then selects the option whose worst grade is highest. Where MEC
asks which option is best on balance, maximin asks which option the most
dissatisfied theory can best live with.

Section~\ref{sec:rules} instantiates one member of each family, MEC for
the expectational approach, the Nash product for bargaining, maximin for
the worst case, and a single-theory \textbf{utilitarian baseline} for a
system that ignores moral uncertainty altogether. Credences are held equal
($w_i = 1/3$) in the main analysis, so differences between rules come
from the aggregation function rather than from a prior ranking of the
traditions. A credence sensitivity analysis later relaxes this assumption
(Section~\ref{sec:results-credences}), and Table~\ref{tab:rules} offers a
concrete mathematical summary of the four rules as they are applied in the
pipeline.

\section{Related Work}
\label{sec:related}
\textbf{RESUME PROOF READING HERE}
Computational treatments of moral uncertainty are recent and sparse.
\citet{ecoffet2021} give the first reinforcement-learning implementation,
encoding rival theories as reward functions in gridworlds and comparing
expectation-style aggregation with voting. They identify intertheoretic
comparison as the immediate practical obstacle. \citet{tennant2024}
fine-tune language-model agents on intrinsic rewards that encode a
utilitarian or a deontological objective, one tradition at a time. More
recent work proposes language-model pipelines that operationalise
moral-uncertainty frameworks for decision support \cite{rosello2025}, or
studies how uncertainty about values surfaces in human and model judgments
\cite{kwon2025}. A parallel line aggregates across people rather than
theories. \citet{conitzer2024} argue that social choice should govern the
aggregation of diverse human feedback in alignment, and the Moral Machine
project distilled
millions of verdicts on autonomous-vehicle dilemmas into a voting-based
decision system \cite{awad2018,noothigattu2018}. None of this work compares
the philosophical aggregation rules as selection criteria over
natural-language advice, which is the setting studied here.

The alignment literature has begun to treat value disagreement as a
problem in its own right. \citet{gabriel2020} distinguishes the technical
problem of fitting a system to given values from the normative problem of
deciding whose values those should be. \citet{sorensen2024} argue for
\emph{pluralistic alignment}, in which a system represents several
reasonable value systems rather than one aggregate. Moral uncertainty is
one formally specified way of doing that, with explicit credences and an
explicit rule in place of an unexamined average. The inclusion of Ubuntu
follows from the same motive. African ethical traditions are nearly absent
from alignment research, although \citet{mhlambi2020} argues that Ubuntu's
relational conception of personhood bears directly on AI governance, and
\citeauthor{metz2007}'s reconstruction of Ubuntu as a moral theory
\cite{metz2007} is precise enough to operationalise. It
enters this pipeline as a first-class framework, with credence equal to
the two Western traditions.

On the evaluation side, ETHICS tests whether models can classify right and
wrong under utilitarian, deontological, virtue, and justice framings
\cite{hendrycks2020}, and Delphi trains a descriptive model of commonsense
moral judgment \cite{jiang2021}. Both score isolated actions rather than
select among competing responses. The dilemmas used here come from
DailyDilemmas, a corpus of everyday quandaries annotated with the values
each side implicates \cite{dailydilemmas2024}; those annotations drive the
filter of Section~\ref{sec:dataset}. Because the judges are language models,
they inherit that method's known biases, verbosity among them
\cite{zheng2023}, and are not taken on trust: the axiom harness of
Section~\ref{sec:axiom-validation} tests each judge directionally, on pairs
length- and register-matched so those biases have little to grip.

The philosophical literature, then, supplies competing rules but no
evidence on whether the choice among them changes the behaviour of a
language-model system. The machine-ethics literature either aggregates
across people or encodes a single tradition, and validates its judges
thinly if at all. This dissertation is the missing experiment, with the
philosophical rules implemented end to end, the judges tested
behaviourally before use, a non-Western tradition included on equal terms,
and divergence measured at the level that matters, the recommended action.

\section{Methodology}
\label{sec:method}

\subsection{Overview}

The pipeline has four stages. First, $500$ dilemmas are drawn from
DailyDilemmas and filtered toward cases where the frameworks are predicted
to disagree (Section~\ref{sec:dataset}). Second, a local
\texttt{llama3.1:8b} model produces sixteen candidate responses per dilemma
under a single neutral prompt, so diversity comes from sampling rather than
framework-flavoured instructions (Section~\ref{sec:generation}). Third,
three prompted \texttt{gpt-4o-mini} judges score every response on a
$0$--$10$ scale, after an axiom harness has checked that each discriminates
in the direction its theory requires (Section~\ref{sec:axiom-validation}).
Fourth, the four selection rules are applied to the globally $z$-scored
ratings, with disagreement measured over the chosen text and the
recommended action (Sections~\ref{sec:aggregation}
and~\ref{sec:divergence}). Credences are equal throughout, and only the
aggregation function varies.

\subsection{Dilemma dataset}
\label{sec:dataset}

DailyDilemmas is a corpus of everyday moral quandaries, each posed as a
binary choice between a named \texttt{TO\_DO} action and a named
\texttt{NOT\_TO\_DO} action, and annotated with the values each side
implicates \cite{dailydilemmas2024}. Those annotations are a keyword
proxy, not a judgment by the later ethical judges. Each annotated value is
mapped onto utilitarian, deontological, or Ubuntu keyword sets. Matches
are counted on each side, and a dilemma is kept only when the three
frameworks split, with at least one preferring \texttt{TO\_DO} and at
least one preferring \texttt{NOT\_TO\_DO}. The kept set is ranked by how
evenly the frameworks pull (\emph{balance}) and by how much annotated
evidence backs the leans (\emph{confidence}), and the top $500$ are
retained.

The filter is a candidate-selection heuristic. Over half of the distinct
values in the source never match a keyword and are dropped, substring
matching misclassifies some antonyms, and the mapping itself is a
judgment. The $500$ dilemmas are therefore a pool expected to be contested,
not a set on which the judges are known to disagree. Whether they do is
settled only by the scores, and part of the observed agreement is an
artefact of the filter admitting dilemmas that proved less torn than the
annotations suggested (Section~\ref{sec:results-read}).

\subsection{Response Generation}
\label{sec:generation}
Sixteen candidate responses per dilemma are sampled from a local
\texttt{llama3.1:8b} model with no framework-flavoured system prompts.
Priming the generator with a moral framework would manufacture the very
disagreement the experiment sets out to detect, so the frameworks enter
only at the judging stage and diversity comes from decoding alone:
temperature $1.1$ with nucleus sampling at $p = 0.95$. The temperature sits
deliberately above the conventional range, because the rules can only be
distinguished by a pool that spreads across the options, and a conservative
setting would return sixteen near-identical answers. The nucleus cap keeps
that heat from degenerating into incoherence. Generation is therefore hot
while scoring is cold (temperature $0.1$,
Section~\ref{sec:axiom-validation}), widening the pool while keeping the
measuring instrument stable, and the homogeneity that nonetheless remains
(Section~\ref{sec:results-read}) is a finding about the generator rather
than an artefact of cautious decoding.

\textbf{A recurring artefact of this stage is refusal}. On a subset of sensitive dilemmas (abortion, sexual addiction,
wishing harm on another person), a substantial fraction of the sixteen
candidates decline to recommend an action, typically with a short ``I
can't\ldots'' disclaimer. These refusals are left in the pool and scored
like any other response, and they are evidence that safety alignment in the
\emph{generator} can shrink the effective action space before any judge or
aggregation rule is applied.

\subsection{Reward Models}
\label{sec:judges}

The three judges are the same underlying model, \texttt{gpt-4o-mini},
steered by a system prompt. Each prompt states the theory's criterion of
rightness, supplies anchors for the full $0$--$10$ range, and asks for a
single integer and nothing else. The utilitarian prompt scores by total
wellbeing and expected outcomes, and the deontological prompt by respect
for persons, perfect duties, and the prohibition on using someone merely
as a means. The Ubuntu prompt scores by harmony, solidarity, and the
quality of the relationships a response would leave behind. The same model is used for all
three so that differences in score can be attributed to the prompt rather
than to differences in base capability.

\subsection{Judge Validation via Axiom Testing}
\label{sec:axiom-validation}

A prompt is not a proof of fidelity. Before the scores can be trusted as a
proxy for utilitarian, deontological, or Ubuntu evaluation, it must be shown
that each judge discriminates between responses in the direction its
assigned theory predicts. Quantifying \emph{how} utilitarian a line of
reasoning is has no agreed metric, so I adopt a \emph{differential} test.
Each theory is operationalised as a small set of defining axioms drawn from
the normative-ethics literature, and the judge must prefer the response that
better satisfies the relevant axiom.

\subsubsection{Operationalising the theories as axioms}
Each framework is decomposed into five operational axioms, each capturing one
central commitment of the theory. These are not offered as canonical
axiomatisations (none of these traditions comes as a settled list of
propositions) but as a defensible operationalisation, grounded in the
literature, of the commitments each judge is meant to track.
Table~\ref{tab:axioms} summarises all fifteen, with the commitment each targets
and the cross-framework contrasts they predict.

\paragraph{Utilitarianism.} Classical utilitarianism
\cite{bentham1789,mill1863} holds that the right action is whichever produces
the most good overall. The independent commitments that structure this claim,
articulated most rigorously by \citet{sidgwick1907}, supply the five
axioms Ut1--Ut5, namely \emph{aggregation} (welfare is summed across all affected
persons, so helping more people is better), \emph{consequentialism} (only
outcomes bear on rightness, not processes), \emph{maximisation under
uncertainty} (expected welfare governs risky choices), the absence of
\emph{side-constraints} (sacrificing one to benefit many is permitted when it
maximises the total), and \emph{welfarism} (only wellbeing is intrinsically
valuable).\footnote{Impartiality receives no axiom of its own. If one option
produces more welfare, Ut1 already tests it, and if two options produce equal
welfare and differ only in \emph{who} is helped, impartiality counts them as
tied and gives no direction to test. It is covered indirectly through the
equal weighting of persons in Ut1 and Ut5.}

\paragraph{Deontology.} Kantian deontology \cite{kant1785} judges an action
by whether its underlying maxim conforms to the moral law, expressed in the
\emph{Formula of Universal Law} (could the maxim be consistently willed as
universal? 4:421)\footnote{References to Kant give the Academy-edition volume
and page.} and the \emph{Formula of Humanity} (persons must always be
treated as ends, never merely as means, 4:429), with \emph{perfect duties}
(strict prohibitions, such as on false promises) distinguished from
\emph{imperfect} ones. Axioms De1--De5 target respect for rational agency,
the means--ends prohibition, the priority of duty over consequences,
universalisability, and the perfect duty of honesty, the last drawing on
Kant's argument that truthfulness is unconditional even where a lie would
prevent harm \cite{kant1797}.

\paragraph{Ubuntu.} Ubuntu locates moral value in the quality of
relationships between people; \citeauthor{mbiti1969}'s ``I am because we
are, and since we are, therefore I am'' \citep[p.106]{mbiti1969} makes
personhood constitutively communal. \citeauthor{metz2007} reconstructs the
tradition as a moral theory, ``an action is right just insofar as it
produces harmony and reduces discord'' \citep[p.338]{metz2007}, where
harmony combines \emph{identity}, the sharing of a way of life, with
\emph{solidarity}, acting for others' good out of genuine good will
\cite{metz2011}. \citeauthor{gyekye1997}'s moderate communitarianism keeps
individual agency exercised within communal life \cite{gyekye1997}, and
\citet{tutu1999} extends the tradition to justice, prioritising the
restoration of harmony over retribution. Ubuntu earns
its place in this study because it generates genuine disagreement with both
Western frameworks wherever they treat the people affected as unconnected
individuals. Axioms Ub1--Ub5 target solidarity, relational awareness,
restorative justice, communal harmony, and relational personhood.

\begin{table*}[t]
\centering
\small
\caption{Fifteen validation axioms, grouped by framework. Each axiom states the property the higher-scoring response should satisfy. Cross-framework contrast indicates where other judges are predicted to disagree. The diagonal column reports own-framework pass rate across the three dilemmas per axiom in a representative run; five-run means are reported in the text, and Ut3 is marginal (see Section~\ref{sec:axiom-results}).}
\label{tab:axioms}
\vspace{4pt}
\begin{tabular}{@{}l p{6.9cm} p{3.6cm} p{2.2cm} c@{}}
\toprule
\textbf{ID} & \textbf{Axiom (higher-scoring response should\ldots)} &
\textbf{Commitment} & \textbf{Cross-fw contrast} & \textbf{Diag.} \\
\midrule
\multicolumn{5}{@{}l}{\textit{Utilitarianism}}\\
Ut1 & help more people                                                  & Aggregation (sum-ranking)         & none            & 3/3 \\
Ut2 & secure the better outcome even through worse process              & Consequentialism                  & Deont.          & 3/3 \\
Ut3 & maximise \emph{expected} wellbeing, even against a certain option & Maximisation under uncertainty    & none            & 2/3 \\
Ut4 & sacrifice one to benefit many                                     & No side-constraints on sacrifice  & Deont.          & 3/3 \\
Ut5 & treat only wellbeing as intrinsically valuable                    & Welfarism                         & none            & 3/3 \\
\addlinespace
\multicolumn{5}{@{}l}{\textit{Deontology}}\\
De1 & respect a competent person's decision rather than override it for their own good & Respect for rational agency       & Util.           & 3/3 \\
De2 & refuse to use a person merely as a means, despite a good outcome  & Formula of Humanity (means--ends) & Util.           & 3/3 \\
De3 & keep a binding duty rather than break it for better consequences  & Priority of duty over consequences& Util.           & 3/3 \\
De4 & refuse an exception for oneself or favoured persons that could not be willed universally & Formula of Universal Law          & none            & 3/3 \\
De5 & tell the truth rather than deceive beneficially                   & Perfect duty of honesty           & Util.           & 3/3 \\
\addlinespace
\multicolumn{5}{@{}l}{\textit{Ubuntu}}\\
Ub1 & act from concern for another's good, with outcomes held equal     & Solidarity                        & none            & 3/3 \\
Ub2 & weigh impact on relationships, not just practicality              & Relational impact                 & none            & 3/3 \\
Ub3 & restore harmony rather than assign blame and punishment           & Restorative over retributive      & Deont.          & 3/3 \\
Ub4 & strengthen communal bonds rather than weaken them                 & Communal harmony                  & Util.           & 3/3 \\
Ub5 & recognise a person as relationally embedded, not atomistic        & Relational personhood             & Deont.          & 3/3 \\
\bottomrule
\end{tabular}
\end{table*}

\subsubsection{Test construction}
For each axiom I author three dilemmas (Ut1a, Ut1b, Ut1c, \ldots, Ub5c), $45$
in total. Every dilemma presents a pair of candidate responses constructed to
differ \emph{only} along the dimension the axiom isolates, with two confounds
controlled explicitly. Pairs are \emph{length-matched} (character-count ratio
$<1.3$, median $1.05$) against verbosity bias, and \emph{register-matched}
(parallel syntax, no one-sided hedging or filler) against a fluency bias that
rewards the more polished response irrespective of content.

Each dilemma carries a predicted direction of preference (\texttt{a > b},
\texttt{b > a}, or null for no prediction) for \emph{all three} judges, not
only the judge whose framework the axiom belongs to. In the
communal-harmony dilemma Ub4b, for instance, a company can move its
meetings online or keep gathering in person. The utilitarian judge is
predicted to prefer the efficiency of moving online, the Ubuntu judge to
prefer the response that keeps the ties of face-to-face contact, and the
deontological judge, indifferent to the dimension varied, carries a null.
The framework-preferred answer alternates between the A and B slots across
dilemmas.

\subsubsection{Scoring and pass criterion}
Crucially, the axiom responses are graded by the \emph{same} judge functions used
in the main pipeline. Each response is presented to the judge \emph{together with
the dilemma it addresses} and scored independently on a $0$--$10$ scale. The judge
never sees the two candidate responses as a pair (so there is no presentation
order to counterbalance), but it does see the scenario that gives a response its
meaning. A predicted cell \emph{passes} when the sign of the score
difference matches the prediction, and an exact tie fails a directional prediction.
A tolerance parameter (set to $0$ here) can discount one-point scoring noise
if raised. Judges score at temperature $0.1$, and owing to the mild
stochasticity the full harness was run five times, with results reported
as averages over those runs.

\subsubsection{Results}
\label{sec:axiom-results}
After the prompt iteration detailed in Section~\ref{sec:held-out}, the
results form a $3 \times 3$ grid of judges against axiom frameworks
(Table~\ref{tab:grid}). The \textbf{diagonal is the validation target}: each
judge on its own framework's axioms. Off-diagonal cells encode predictions
about where the frameworks should \emph{disagree}, so a failed off-diagonal
cell is a failed prediction about cross-framework contrast, not evidence
that a judge misapplies its own theory.

Across the $75$ cells carrying a non-null prediction, the judges matched the
predicted direction on $86.9\%$ on average (per run: 88, 87, 87, 85,
$88\%$). On the diagonal they are near-perfect, at a mean of $99.1\%$; the
deontological and Ubuntu judges pass all fifteen of their own cells in every
run.

\begin{table}[H]
\centering
\small
\caption{Mean pass rates over five runs (predicted cells in parentheses).
Diagonal entries in bold, and --- marks a null prediction.}
\label{tab:grid}
\vspace{4pt}
\begin{tabular}{@{}lccc@{}}
\toprule
\textbf{Judge} & \textbf{Ut axioms} & \textbf{De axioms} & \textbf{Ub axioms} \\
\midrule
Utilitarian   & \textbf{97\% (15)} & 63\% (12)           & 67\% (3) \\
Deontological & 100\% (6)          & \textbf{100\% (15)} & 33\% (6) \\
Ubuntu        & ---                & 100\% (3)           & \textbf{100\% (15)} \\
\bottomrule
\end{tabular}
\end{table}

The only diagonal instability is Ut3a, a risk-neutrality case pitting a
one-in-five chance of saving $100$ lives (expected value $20$) against saving
$5$ for certain; the utilitarian judge passes in three of five runs, scoring
the options within a point of one another. This reflects a residual,
alignment-instilled caution toward gambles that will probably yield nothing
rather than an inability to reason as a utilitarian: the same judge takes the
gamble decisively when the upside is very large (Ut3c, a one-in-ten chance of
helping $100{,}000$), so it weighs expected value but discounts risky options
near indifference.

The remaining misses are all off-diagonal and, aside from De2a under the
utilitarian judge (3/5), stable across runs. They fall into two classes:
philosophically contestable predictions (the deontological judge preferred
restorative to retributive responses on all three Ub3 dilemmas; the
utilitarian judge missed two of three De1 respect-for-agency cells) and
near-ties, where both responses score within a point in every run so no
direction emerges (Ub4c; Ub5b under the deontological judge). The weakest
entry in Table~\ref{tab:grid}, the deontological judge on Ubuntu axioms
($33\%$), is exactly these two classes alongside a perfect diagonal.

One substantive finding deserves emphasis. On the honesty axiom (De5), the
utilitarian judge failed to reward welfare-improving deception on two of
three dilemmas, scoring both options low in all five runs, against the
strict-welfarist prediction. Alignment training appears to instil an honesty
guardrail that persists even when the model is instructed to reason as a
strict consequentialist, a limitation worth bearing in mind whenever aligned
LLMs are used as ethical evaluators.

\subsubsection{Held-out validation}
\label{sec:held-out}
Two judge prompts were revised in response to observed axiom behaviour, which
raises a circularity concern: a clause added after seeing a failure might fit
those dilemmas rather than reflect a general capability. The deontological
judge gained a clause making truthfulness a perfect duty that outranks a
better outcome secured by lying; the utilitarian judge lost a clause (``a
larger but unlikely benefit does not automatically outrank a smaller but
more certain one'') that encoded risk-\emph{aversion} rather than classical
risk-neutrality, a theory-faithfulness correction that improved Ut3 from a
reliable miss to the marginal pass above.

To test for overfitting, I froze the prompts, authored six fresh dilemmas
of the same construction (three honesty, three risk-neutrality), and ran
five times. The honesty clause generalised cleanly, all three new dilemmas
passing in every run; the risk-neutrality guidance generalised only
directionally, the three new cases passing in five, three, and four runs
with scores again compressed near indifference. Held-out performance
overall was $90\%$: the honesty edit reflects a genuine disposition, the
risk edit correct-but-marginal behaviour consistent with the residual
risk-aversion above.

\paragraph{Iteration and limitations.} Two further prompt improvements were
axiom-independent: vague ``be harsh'' instructions became the anchored
$0$--$10$ rubrics after early runs showed severe score compression, and
responses are scored in the context of their dilemma rather than in
isolation. Caveats remain. Each axiom rests on only three dilemmas, and the
harness validates \emph{directional discrimination on constructed,
deliberately clean cases}, not correct reasoning on the messier dilemmas of
the main study. Cross-framework coverage is also uneven: the Ubuntu judge is
predicted to contrast with the others on only one non-Ubuntu axiom (De2), so
the harness gives little evidence that it \emph{diverges} from the other
judges as the theory predicts, leaving open the possibility that it is a
generically pro-communal scorer rather than a distinctively Ubuntu one.

\subsection{Aggregation}
\label{sec:aggregation}

\subsubsection{The aggregation rules}
\label{sec:rules}

For a dilemma with candidate responses $R$, let $c_i(R)$ be the
choiceworthiness of $R$ under framework $i \in \{\text{Ut}, \text{De},
\text{Ub}\}$, with equal credences $w_i = 1/3$ in the main analysis so that
differences between rules come from the aggregation function rather than
from a prior ranking of the traditions (a sensitivity analysis relaxes
this; Section~\ref{sec:results-credences}). Each rule selects
\[
  R^{\star} = \arg\max_{R}\, f\bigl(c_1(R),\, c_2(R),\, c_3(R)\bigr),
\]
so the four rules differ only in $f$ (Table~\ref{tab:rules}); what is
substituted for $c_i$ is settled in Sections~\ref{sec:intertheoretic}
and~\ref{sec:zscore}. The \textbf{baseline} consults the utilitarian judge
alone, standing in for a system aligned to one tradition and giving the
comparison point for what aggregation buys. \textbf{Expected
choiceworthiness} averages the three scores \cite{macaskill2020}, so a large
enough gain for one theory can outvote the others. \textbf{Maximin} selects
the response whose worst-treated framework fares best, refusing the trade EC
allows; it depends on levels, not only rankings. The \textbf{Nash moral
parliament} maximises the product of each framework's surplus $S_i$ above a
disagreement point \cite{nash1950}, rewarding balance: a response scoring
near zero for any framework has its product destroyed however well the
others rate it. (The weighted Nash solution $\prod_i S_i^{w_i}$ reduces to
the plain product under equal credences, since a common power preserves the
ordering.) Raw-score ties are common and survive normalisation; the
convention for handling them is given in Section~\ref{sec:ties}.

\begin{table}[htbp]
\centering
\small
\caption{The four selection rules. $c_i$ is the normalised choiceworthiness
under framework $i$, and $S_i$ the Nash surplus defined in
Section~\ref{sec:zscore}.}
\label{tab:rules}
\vspace{2pt}
\begin{tabular}{@{}l c p{3.0cm}@{}}
\toprule
\textbf{Rule} & \textbf{Objective $f$} & \textbf{Character} \\
\midrule
Baseline & $c_{\text{Ut}}$ & Single tradition; disagreement ignored \\
EC       & $\sum_i w_i c_i$ & Averages; permits trade-offs \\
Maximin  & $\min_i c_i$ & Protects the worst-served framework \\
Nash     & $\prod_i S_i$ & Rewards balance; punishes near-zero scores \\
\bottomrule
\end{tabular}
\end{table}

\subsubsection{The intertheoretic value problem}
\label{sec:intertheoretic}

Raw $0$--$10$ scores cannot be aggregated as if they were already comparable.
A judge that spreads its scores widely (say $2$--$8$) will swing EC, Maximin,
and Nash more than a judge stuck in a narrow band ($7$--$9$), even under
equal credences. This is the \textbf{intertheoretic value problem} of
Section~\ref{sec:background} in statistical form \cite{macaskill2020}:
plugging unnormalised scores into a shared formula lets scale artefacts, not
only moral content, decide the outcome. The axiom harness never faces this,
since it checks ordering \emph{within} a single judge, but normalisation is
required once judges are combined. Scores are therefore normalised so that
each framework influences the choice by how it ranks the candidates, not by
how much of the scale it happens to use.

\subsubsection{Global z-scoring}
\label{sec:zscore}

Each judge $j$'s pooled raw scores over every dilemma and candidate yield a
mean $\mu_j$ and standard deviation $\sigma_j$, and each raw score $x$ is
replaced with $z = (x - \mu_j)/\sigma_j$. Every framework is then centred at
zero with unit spread, so a one-point lead no longer counts for more under a
wide-ranging judge than under a compressed one.

The scope is deliberately \emph{global} rather than per-dilemma. Per-dilemma
$z$-scoring divides by the standard deviation of that dilemma's sixteen
candidates alone; when a judge scores them all within a narrow band that
deviation is tiny, a one-point gap is stretched into a large one, and a judge
that rated every candidate between $8$ and $9$ is reported as decisively
preferring one of them. Global $\mu_j$ and $\sigma_j$ measure each gap
against how much the judge varies across the whole corpus, so genuine
near-indifference stays small.

EC and Maximin use these $z$-scores directly. Nash cannot, since products of
negative values are not meaningful as choiceworthiness; it needs non-negative
gains above a disagreement point, fixed a priori at a raw score of $0$
(absence of moral worth on the rubric), not at the worst score observed in
the sample. Subtracting the $z$-score of raw $0$ from each response's
$z$-score simplifies to the surplus $S_i = X_i / \sigma_i$, where $X_i$ is
the raw score under judge $i$. Because the Nash product is scale-invariant,
dividing by $\sigma_i$ does not change which response is selected relative to
a raw product of the $X_i$; the substantive choice is the disagreement point
at raw $0$.

\subsection{Divergence Measurement}
\label{sec:divergence}

Divergence is not a single quantity. Two rules may select different response
texts that direct the reader toward the same course of action, which is
disagreement over justification and rhetoric, or they may select texts whose
core recommendations are opposed, which is disagreement over what ought to
be done. The first reveals how the choice of rule shapes the reasoning a
system presents; the second is the level at which the research question is
posed. Both are measured and reported separately.

\subsubsection{Recommendation labelling}
\label{sec:labelling}

The source dataset records, for every dilemma, a pair of canonical actions: a
\texttt{TO\_DO} action and a \texttt{NOT\_TO\_DO} action. Each of the sixteen
candidate responses is assigned one of four labels according to the action it
endorses:

\begin{itemize}[leftmargin=*]
  \item \texttt{TO\_DO}: the response points the reader toward the named
  to-do action.
  \item \texttt{NOT\_TO\_DO}: it points toward the named not-to-do action.
  \item \texttt{REFUSAL}: it declines to engage with the dilemma at all
  (typically a short ``I cannot advise on this'' disclaimer, with no
  recommendation attached).
  \item \texttt{OTHER}: a distinct third course, or no discernible lean
  toward either named action.
\end{itemize}

The threshold for the two action labels is deliberately permissive: a
response that recommends an action tentatively, conditionally, or in
softened form is still recorded as endorsing it, since the question is which
way the text would move a reader, not how forcefully it argues. An earlier
revision of the prompt reserved the action labels for confident endorsements
and routed cautious answers into a residual ``hedge'' category, which
conflated hesitancy with refusal and split near-identical responses across
categories; \texttt{REFUSAL} is accordingly restricted to responses that
offer no recommendation of any kind.

Labels are assigned \emph{post hoc} by a separate model
(\texttt{gpt-4o-mini}) rather than emitted by the generator alongside each
response: a self-reported label is a claim about intent rather than an
independent reading of the text, and requiring the generator to commit to a
label while writing risks the label shaping the prose. Post-hoc annotation
keeps the labels independent of generation at the cost of a second source of
error, so each label is stored with the short phrase in the response that
carries the recommendation, allowing the labelling to be audited without
re-reading the corpus.

Two residual weaknesses are acknowledged at the outset. \texttt{OTHER} is a
residual category, so a shared \texttt{OTHER} is weaker evidence of
agreement than a shared action label (two responses may endorse quite
different third courses); after labelling it proves small, at $63$ responses
(Section~\ref{sec:results-labels}). And the labels have not been validated
against human ground truth: automatic consistency checks are reported in
Section~\ref{sec:results-labels}, and a hand-labelled gold sample is left to
future work.

\subsubsection{Set-valued winners}
\label{sec:ties}

On a $0$--$10$ integer rubric, distinct responses frequently receive
identical score patterns, and because global $z$-scoring is an affine
transform of each judge's raw scores, such ties survive normalisation
intact. Several candidates may therefore share the top value under a given
rule. Rather than breaking ties arbitrarily, each rule's selection is
treated as a \emph{winner set}: all candidates whose value lies within
$10^{-9}$ of the maximum, a tolerance that only absorbs floating-point noise
and is orders of magnitude smaller than any genuine gap in the data. A
random tie-break would have made the results depend on an arbitrary seed;
under the set-valued convention, every reported disagreement is one that no
choice of tie-break could remove.

\subsubsection{Agreement at two levels}
\label{sec:two-levels}

Let $S_A$ and $S_B$ denote the winner sets of rules $A$ and $B$ on a given
dilemma, and let $\ell(\cdot)$ denote the recommendation label of a response.

\paragraph{Response level.}
Rules $A$ and $B$ \emph{can agree} on a response when their winner sets
overlap, $S_A \cap S_B \neq \emptyset$, and \emph{diverge} when the sets are
disjoint. In the disjoint case no tie-breaking convention, however
favourable, could produce a common winner, so the disjoint count is a lower
bound on disagreement. A corresponding upper bound, used only as a
robustness check, is the probability that two uniform draws from $S_A$ and
$S_B$ select different responses,
$1 - |S_A \cap S_B| / (|S_A| \cdot |S_B|)$.

\paragraph{Recommendation level.}
Rules $A$ and $B$ can agree on an \emph{action} when some label appears in
both winner sets,
$\{\ell(r) : r \in S_A\} \cap \{\ell(r) : r \in S_B\} \neq \emptyset$, and
diverge when the label sets share nothing, as when one winner set is
entirely \texttt{TO\_DO} and the other entirely \texttt{NOT\_TO\_DO}. This
criterion answers whether the rules would direct a user toward different
courses of action.

Both definitions extend to all four rules at once: the rules are unanimous
at a level when the intersection over all four winner sets (respectively,
label sets) is non-empty. Response-level unanimity entails
recommendation-level unanimity, since a shared response carries a shared
label; the converse fails, and the gap between the two is the quantity of
interest in Section~\ref{sec:results-two-levels}. Two kinds of nominal
agreement are not treated as agreement on a moral position and are counted
separately: a shared \texttt{REFUSAL} records only that the generator
declined to answer, and a shared \texttt{OTHER} only that both winners
depart from the binary, not that they depart in the same direction.

\section{Results}
\label{sec:results}

The pipeline is applied to all $500$ dilemmas: sixteen candidate responses
per dilemma, scored by the three judges, normalised as in
Section~\ref{sec:zscore}, and ranked by the four rules of
Table~\ref{tab:rules}. The results are presented in five stages: the
composition of the labelled pool, the prevalence of ties, agreement at the
response level, agreement at the recommendation level, and the sensitivity
of that agreement to the credences.

\subsection{The labelled pool}
\label{sec:results-labels}

Of the $8{,}000$ candidate responses, $3{,}195$ are labelled
\texttt{TO\_DO}, $4{,}074$ \texttt{NOT\_TO\_DO}, $668$ \texttt{REFUSAL}, and
$63$ \texttt{OTHER}. Two features of this distribution bear on what follows.
Refusals are substantial and structurally distinct: the median refusal runs
to $25$ words against $308$ for other responses, and all but one of the
$668$ open with an explicit decline. This is the generator-side safety
behaviour noted in Section~\ref{sec:generation}, and the aggregation rules
do sometimes select such responses (refusals occupy forty-six winning slots
across the four rules), so a safety refusal can outrank every substantive
recommendation available for a dilemma. \texttt{OTHER}, by contrast, is rare
($0.8\%$), a genuine residual rather than a repository for cautious answers,
which limits the exposure identified in Section~\ref{sec:labelling}.

Because the labels were produced by a second language model, two automatic
consistency checks were run. Of the $7{,}659$ labels carrying a supporting
quote, $99.1\%$ of the quotes appear verbatim in the response they were
drawn from once punctuation and markdown are normalised, the remainder being
light paraphrases, which is evidence that the labeller reads the text rather
than inventing support. Responses that open with a refusal phrase yet
received an action label prove, on inspection, to follow the ``disclaimer,
then advice'' pattern the prompt instructs the labeller to judge on the
advice. These checks establish internal consistency, not accuracy:
recommendation-level figures carry that caveat until a hand-labelled
sample exists.

\subsection{Ties}
\label{sec:results-ties}

Ties at the top are pervasive and markedly uneven across methods
(Table~\ref{tab:ties}). The utilitarian baseline is indifferent among two or
more candidates on $383$ of $500$ dilemmas ($77\%$), with a mean winner set
of seven when it ties and $26$ dilemmas where its winner set is the entire
pool. EC and Nash tie on roughly a third of dilemmas, Maximin on just over
half.

The pattern reflects the scoring scale more than moral indifference. Across
the sixteen candidates for a dilemma, a judge draws on only $3.6$ distinct
utilitarian values on average ($3.6$ deontological, $4.0$ Ubuntu), and the
candidates realise a mean of $9.4$ distinct score triples, so $60\%$ of
responses share their exact triple with a competitor for the same dilemma;
identical triples remain tied after $z$-scoring, which is affine. The
distributions are also top-heavy (the deontological judge places $58\%$ of
its scores in the $8$--$10$ band, the other judges roughly $41\%$), the
ceiling effect anticipated during judge design. The baseline inherits this
coarseness directly, consulting one coordinate, whereas the aggregation
rules combine three and so resolve some ties the baseline cannot. The
practical consequence is the convention of Section~\ref{sec:ties}:
disagreement is computed on winner sets, since a large baseline tie-set that
happens not to contain the EC winner is response-level divergence even when
every text in both sets recommends the same action. Exactly this possibility
motivates the two-level analysis.

\begin{table}[htbp]
\centering
\small
\caption{How often each rule is indifferent among two or more top-scoring
responses ($n = 500$ dilemmas). ``Complete'' means the winner set is all
sixteen candidates.}
\label{tab:ties}
\vspace{2pt}
\begin{tabular}{@{}l r r r r@{}}
\toprule
\textbf{Rule} & \textbf{Ties} & \textbf{Rate} & \textbf{Mean size} & \textbf{Complete} \\
\midrule
EC        & 166 & $33\%$ & $3.3$ & 0 \\
Maximin   & 280 & $56\%$ & $4.9$ & 6 \\
Nash      & 180 & $36\%$ & $3.3$ & 0 \\
Baseline  & 383 & $77\%$ & $7.1$ & 26 \\
\bottomrule
\end{tabular}
\end{table}

\subsection{Two levels of agreement}
\label{sec:results-two-levels}

Table~\ref{tab:two-levels} presents the central result. On $339$ of $500$
dilemmas ($68\%$), some single response lies in every rule's winner set, so
the four rules can agree on both the text and the action. On a further $105$
($21\%$) no common text exists, yet every winner set still contains the same
core recommendation. Only on the remaining $56$ ($11\%$) do the rules
diverge on the action itself.

\begin{table}[htbp]
\centering
\small
\caption{Agreement of all four rules, at two levels ($n = 500$). Sharing
a response implies sharing its label, so the missing cell is empty by
construction.}
\label{tab:two-levels}
\vspace{2pt}
\begin{tabular}{@{}l r r@{}}
\toprule
 & \textbf{Count} & \textbf{Share} \\
\midrule
Same response and same recommendation     & 339 & $67.8\%$ \\
Different response, same recommendation   & 105 & $21.0\%$ \\
Different recommendation                  & 56  & $11.2\%$ \\
\bottomrule
\end{tabular}
\end{table}

The figure of $56$ answers the research question as it concerns \emph{what
to do}, and it is a lower bound: a dilemma counts as action-level divergence
only when the four label sets are entirely disjoint. Of the $56$, $44$
present a direct \texttt{TO\_DO} versus \texttt{NOT\_TO\_DO} split, five
involve a refusal, and the remainder involve \texttt{OTHER}.

The $105$ dilemmas in the middle row are those where the rules prefer
different texts that point the same way. In a minority the difference is a
genuine difference of justification, as when one winner counsels ``confess,
but carefully'' where another counsels simply telling the truth, both
\texttt{TO\_DO}. In rather more, the winner sets are non-overlapping subsets
of near-paraphrases, typically because the baseline is indifferent among a
large collection of similar answers from which the aggregation rules select
a different subset. Twenty-four of the $105$ are unique-winner splits, the
strongest evidence that the choice of rule can change which justification is
presented even when the action is unchanged. The raw response-level figure
of $161/500$ ($32\%$) therefore overstates practical disagreement
considerably: most of it is rhetorical, and some of the rhetoric is
paraphrase.

Of the $444$ recommendation-level agreements, $221$ fall on
\texttt{NOT\_TO\_DO}, $212$ on \texttt{TO\_DO}, nine on \texttt{REFUSAL},
and two on \texttt{OTHER}. The nine refusal-agreements record only that the
generator declined to answer; they are retained in the $444$ but excluded
from any claim that the rules share a moral position.

\subsection{The structure of the disagreement}
\label{sec:results-pairwise}

Table~\ref{tab:pairwise} reports every pair of rules at both levels, and the
pattern is consistent. EC and Nash almost never select different actions
($4$ of $500$ dilemmas, $0.8\%$); Maximin departs from either more often
but still in low single digits; and nearly all action-level disagreement
lies between the three uncertainty rules and the utilitarian baseline, a
contrast that is itself modest at $38$--$44$ dilemmas per pair (roughly
$8\%$). Consistent with this, of the $56$ dilemmas on which the four rules
cannot share a recommendation, $31$ have the three uncertainty rules
agreeing on an action with the baseline apart; the remaining $25$ are splits
among the uncertainty rules themselves.

\begin{table}[htbp]
\centering
\small
\caption{Pairwise divergence ($n = 500$). Response-level: winner sets
are disjoint. Recommendation-level: label sets are disjoint.}
\label{tab:pairwise}
\vspace{2pt}
\begin{tabular}{@{}l r r@{}}
\toprule
\textbf{Pair} & \textbf{Response} & \textbf{Recommendation} \\
\midrule
EC vs Nash            & $7.4\%$  & $0.8\%$ \\
Maximin vs Nash       & $11.0\%$ & $3.8\%$ \\
EC vs Maximin         & $20.6\%$ & $5.0\%$ \\
EC vs Baseline        & $16.6\%$ & $7.8\%$ \\
Nash vs Baseline      & $21.4\%$ & $7.6\%$ \\
Maximin vs Baseline   & $26.6\%$ & $8.8\%$ \\
\bottomrule
\end{tabular}
\end{table}

The forced-choice figures are considerably larger (around $70$--$80\%$
against the baseline, $29\%$ even for EC against Nash) because a large
tie-set makes a uniform draw unlikely to land on the one shared candidate.
They are reported only as an upper bound: that randomness is no part of the
rules, and treating it as a headline would attribute to them a disagreement
that belongs to the tie-break.

\subsection{Credence sensitivity}
\label{sec:results-credences}

To test whether the observed agreement depends on the equal-credence
choice, the aggregation stage was re-run twice with a $2{:}1{:}1$ tilt:
credences $(0.25, 0.5, 0.25)$ favouring deontology, and $(0.25, 0.25,
0.5)$ favouring Ubuntu. Only EC and Nash can respond (EC reweights its
sum; Nash becomes the asymmetric product $\prod_i S_i^{w_i}$); Maximin
protects the worst-served theory whatever its credence and the baseline
consults a single judge, so both are unchanged by construction.

The tilts move the rules considerably at the level of text: the
deontological tilt hands EC a different winner set on $85$ of $500$
dilemmas and Nash on $70$ (the Ubuntu tilt: $49$ and $55$). Almost none
of that movement reaches the action (Table~\ref{tab:credences}).
Action-level disagreement rises from $56$ dilemmas to $61$ under either
tilt, and most dilemmas that lose their shared response merely prefer a
different text pointing the same way. What growth there is concentrates
against the credence-free baseline (EC versus baseline, $7.8\%$ to
$10.4\%$ under the deontological tilt), while EC versus Nash stays near
$1\%$. Doubling a theory's credence changes which justification wins far
more often than what the system recommends; the headline of
Table~\ref{tab:two-levels} is not an artefact of the equal-credence
choice.

\begin{table}[htbp]
\centering
\small
\caption{Agreement of all four rules under three credence settings
($n = 500$). Each tilted run gives the named framework credence $0.5$ and
the other two $0.25$.}
\label{tab:credences}
\vspace{2pt}
\begin{tabular}{@{}l r r r@{}}
\toprule
 & \textbf{Equal} & \textbf{De $0.5$} & \textbf{Ub $0.5$} \\
\midrule
Same response, same recommendation      & 339 & 316 & 323 \\
Different response, same recommendation & 105 & 123 & 116 \\
Different recommendation                & 56  & 61  & 61  \\
\bottomrule
\end{tabular}
\end{table}

\subsection{Interpretation}
\label{sec:results-read}

On the quantity the motivating premise concerns, the recommended action, the
observed effect is modest. Once justification is set aside, EC and Nash are
close to interchangeable, and replacing a utilitarian default with any of
the three uncertainty rules changes the selected action in fewer than one
dilemma in ten, on a pool already filtered toward moral contestedness.

The result remains informative because it locates \emph{where} the machinery
matters: rarely between one uncertainty rule and another, occasionally
against a single-tradition default, and more often as a change of
justification than of action. Three features of the pipeline help explain
why the action-level figure is small, and they belong alongside the result.
First, the candidate pool is homogeneous: responses come from an
$8$B-parameter model under one neutral prompt, so many of the sixteen
candidates are close restatements of the same side, and every rule is
reduced to ranking stylistic variants of one recommendation. Second, the
judges sometimes saturate: where nearly every candidate satisfies the
relevant duty or relational obligation, the deontological and Ubuntu scores
cluster near the top of the scale, EC, Maximin, and Nash become effectively
utilitarian-driven, and agreement there is an artefact of compressed
scores. Third, the keyword filter of Section~\ref{sec:dataset} predicts
where disagreement should arise but cannot guarantee it, so part of the
pool proved less contested than the annotations suggested.

None of this licenses substituting the larger response-level count for the
$56$: that count is real but answers a different question, which text is
preferred, and it is inflated by ties among paraphrases. The defensible
claim is the narrower one Table~\ref{tab:pairwise} supports: modelling moral
uncertainty rarely changes the recommended action relative to other
uncertainty rules, and changes it relative to a utilitarian default in
roughly eight dilemmas in a hundred.

\section{Future Work}

Four limitations point to concrete extensions.

\emph{The measurement scale is coarse}, and the tie rates of
Section~\ref{sec:results-ties} are its most visible consequence: a judge
that separates sixteen candidates into fewer than four distinct values
cannot express fine distinctions, so much of the observed indifference is an
artefact of resolution. Re-scoring on a wider rubric is not the simple fix
it appears: language-model judges cluster on round numbers, so a $0$--$100$
scale would resolve only rounding ties, not near-paraphrase ties, and the
anchors were axiom-validated at their present granularity, so rewriting
them yields a new instrument that must be revalidated before any
comparison. Meanwhile the set-valued convention keeps the coarseness from
inflating the headline.

\emph{The recommendation labels require human validation.} The automatic
checks establish that the labeller is consistent with its own prompt and its
quotes grounded in the text, not that a human would assign the same labels.
A stratified gold sample of several hundred responses, labelled blind to the
model's output, would give per-class precision (most importantly on the
softened-action boundary) and a prevalence-reweighted estimate of corpus
accuracy. Until then, the $56$ action-level disagreements are the output of
an unvalidated classifier atop an otherwise fully specified pipeline.

\emph{A categorical label discards degree.} ``Gradually distance yourself''
and ``cut off contact immediately'' may share a \texttt{NOT\_TO\_DO} tag
while differing sharply in force. A scalar lean on $[-1, +1]$ recorded
alongside the category would let divergence be graded rather than binary,
though it should not be folded into the aggregation scores, which measure
choiceworthiness, not how firmly a response is phrased.

\emph{The homogeneity of the candidate pool is the binding constraint}, as
Section~\ref{sec:results-read} argues. A larger generator, or a procedure
designed to produce genuinely opposed courses of action rather than sixteen
variants of one side, would give the rules materially more to discriminate
between. Were action-level divergence to remain near $11\%$ under such a
pool the negative result would be considerably strengthened, and were it to
rise the present figure would stand revealed as an artefact of sampling
rather than a fact about the rules.

\begin{thebibliography}{99}
\footnotesize
\bibitem[Arrow(1951)]{arrow1951} Arrow, K. J. (1951). \textit{Social Choice and Individual Values}. Yale University Press.
\bibitem[Awad et~al.(2018)]{awad2018} Awad, E., Dsouza, S., Kim, R., Schulz, J., Henrich, J., Shariff, A., Bonnefon, J.-F., \& Rahwan, I. (2018). The Moral Machine experiment. \textit{Nature}, 563(7729), 59--64.
\bibitem[Bentham(1789)]{bentham1789} Bentham, J. (1789). \textit{An Introduction to the Principles of Morals and Legislation}. London: T. Payne \& Son.
\bibitem[Bostrom(2009)]{bostrom2009} Bostrom, N. (2009). Moral uncertainty --- towards a solution? \textit{Overcoming Bias}.
\bibitem[Casper et~al.(2023)]{casper2023} Casper, S., et al. (2023). Open problems and fundamental limitations of reinforcement learning from human feedback. \textit{Transactions on Machine Learning Research}.
\bibitem[Christian(2020)]{christian2020} Christian, B. (2020). \textit{The Alignment Problem: Machine Learning and Human Values}. W. W. Norton \& Company.
\bibitem[Conitzer et~al.(2024)]{conitzer2024} Conitzer, V., et al. (2024). Aggregation problems in machine ethics and AI alignment. \textit{AAAI/ACM Conference on AI, Ethics, and Society}.
\bibitem[Cotton-Barratt et~al.(2020)]{cottonbarratt2020} Cotton-Barratt, O., MacAskill, W., \& Ord, T. (2020). Statistical normalization methods in interpersonal and intertheoretic comparisons. \textit{Journal of Philosophy}, 117(2), 61--95.
\bibitem[Chiu et~al.(2024)]{dailydilemmas2024} Chiu, J., et al. (2024). DailyDilemmas: Revealing value preferences of LLMs with quandaries of daily life. \textit{ICLR 2025}.
\bibitem[Ecoffet and Lehman(2021)]{ecoffet2021} Ecoffet, A. \& Lehman, J. (2021). Reinforcement learning under moral uncertainty. \textit{Proceedings of ICML}.
\bibitem[Gabriel(2020)]{gabriel2020} Gabriel, I. (2020). Artificial intelligence, values, and alignment. \textit{Minds and Machines}, 30(3), 411--437.
\bibitem[Greaves and Cotton-Barratt(2019)]{greaves2019} Greaves, H. \& Cotton-Barratt, O. (2019). A bargaining-theoretic approach to moral uncertainty. \textit{Global Priorities Institute Working Paper} No. 4-2019, University of Oxford.
\bibitem[Gyekye(1997)]{gyekye1997} Gyekye, K. (1997). \textit{Tradition and Modernity: Philosophical Reflections on the African Experience}. Oxford University Press.
\bibitem[Hendrycks et~al.(2020)]{hendrycks2020} Hendrycks, D., et al. (2020). Aligning AI with shared human values. \textit{arXiv:2008.02275}.
\bibitem[Jiang et~al.(2021)]{jiang2021} Jiang, L., Hwang, J. D., Bhagavatula, C., Le Bras, R., Forbes, M., Borchardt, J., Liang, J., Etzioni, O., Sap, M., \& Choi, Y. (2021). Delphi: Towards machine ethics and norms. \textit{arXiv:2110.07574}.
\bibitem[Kant(1785)]{kant1785} Kant, I. (1785). \textit{Groundwork of the Metaphysics of Morals}. (M. Gregor, Trans., 1998). Cambridge University Press.
\bibitem[Kant(1797)]{kant1797} Kant, I. (1797). On a supposed right to lie from philanthropic concerns. \textit{Berlinische Bl\"{a}tter}, 1, 301--314.
\bibitem[Kwon et~al.(2025)]{kwon2025} Kwon, J., et al. (2025). Dropouts in confidence: Moral uncertainty in human-LLM alignment. \textit{arXiv:2511.13290}.
\bibitem[Lockhart(2000)]{lockhart2000} Lockhart, T. (2000). \textit{Moral Uncertainty and Its Consequences}. Oxford University Press.
\bibitem[MacAskill et~al.(2020)]{macaskill2020} MacAskill, W., Bykvist, K., \& Ord, T. (2020). \textit{Moral Uncertainty}. Oxford University Press.
\bibitem[Mbiti(1969)]{mbiti1969} Mbiti, J. S. (1969). \textit{African Religions and Philosophy}. Heinemann.
\bibitem[Metz(2007)]{metz2007} Metz, T. (2007). Toward an African Moral Theory. \textit{Journal of Political Philosophy}, 15(3), 321--341.
\bibitem[Metz(2011)]{metz2011} Metz, T. (2011). Ubuntu as a moral theory and human rights in South Africa. \textit{African Human Rights Law Journal}, 11(2), 532--559.
\bibitem[Mhlambi(2020)]{mhlambi2020} Mhlambi, S. (2020). From rationality to relationality: Ubuntu as an ethical and human rights framework for artificial intelligence governance. \textit{Carr Center Discussion Paper}, Harvard Kennedy School.
\bibitem[Mill(1863)]{mill1863} Mill, J. S. (1863). \textit{Utilitarianism}. London: Parker, Son \& Bourn.
\bibitem[Nash(1950)]{nash1950} Nash, J. (1950). The bargaining problem. \textit{Econometrica}, 18(2), 155--162.
\bibitem[Newberry and Ord(2021)]{newberry2021} Newberry, T. \& Ord, T. (2021). The parliamentary approach to moral uncertainty. \textit{Future of Humanity Institute Technical Report} \#2021-2, University of Oxford.
\bibitem[Noothigattu et~al.(2018)]{noothigattu2018} Noothigattu, R., Gaikwad, S., Awad, E., Dsouza, S., Rahwan, I., Ravikumar, P., \& Procaccia, A. D. (2018). A voting-based system for ethical decision making. \textit{Proceedings of AAAI}.
\bibitem[Ouyang et~al.(2022)]{ouyang2022} Ouyang, L., et al. (2022). Training language models to follow instructions with human feedback. \textit{NeurIPS}.
\bibitem[Rosello et~al.(2025)]{rosello2025} Rosello, P., et al. (2025). Addressing moral uncertainty using large language models for ethical decision-making. \textit{arXiv:2503.05724}.
\bibitem[Sidgwick(1907)]{sidgwick1907} Sidgwick, H. (1907). \textit{The Methods of Ethics} (7th ed.). London: Macmillan.
\bibitem[Sorensen et~al.(2024)]{sorensen2024} Sorensen, T., Moore, J., Fisher, J., Gordon, M., Mireshghallah, N., Rytting, C. M., Ye, A., Jiang, L., Lu, X., Dziri, N., Althoff, T., \& Choi, Y. (2024). A roadmap to pluralistic alignment. \textit{Proceedings of ICML}.
\bibitem[Tennant et~al.(2024)]{tennant2024} Tennant, E., et al. (2024). Moral alignment for LLM agents. \textit{arXiv:2410.01639}.
\bibitem[Tutu(1999)]{tutu1999} Tutu, D. (1999). \textit{No Future Without Forgiveness}. Rider Books.
\bibitem[Zheng et~al.(2023)]{zheng2023} Zheng, L., Chiang, W.-L., Sheng, Y., Zhuang, S., Wu, Z., Zhuang, Y., Lin, Z., Li, Z., Li, D., Xing, E., Zhang, H., Gonzalez, J. E., \& Stoica, I. (2023). Judging LLM-as-a-Judge with MT-Bench and Chatbot Arena. \textit{NeurIPS}.
\end{thebibliography}

\end{document}
